% A könyv borítója és PDF-címoldala (B5, 176 × 250 mm) a Pannon Egyetem arculatával.
% Az EPUB-borító (assets/cover.png) ebből készül: pdftocairo -png -r 230 -singlefile
\documentclass[tikz,border=0pt]{standalone}
\usepackage{pe-tikz}
\usepackage{graphicx}
\usepackage{fontspec}
\newfontfamily\cvblack{Raleway}[Path=../../assets/fonts/, Extension=.otf, UprightFont=*-Black]
\newfontfamily\cvbold{Raleway}[Path=../../assets/fonts/, Extension=.otf, UprightFont=*-Bold]
\newfontfamily\cvserif{IBMPlexSerif}[Path=../../assets/fonts/, Extension=.otf,
  UprightFont=*-Regular, ItalicFont=*-Italic]
\newfontfamily\cvsans{IBMPlexSans}[Path=../../assets/fonts/, Extension=.otf,
  UprightFont=*-Regular, BoldFont=*-SemiBold]
\definecolor{cvblue}{HTML}{285780}
\definecolor{cvmid}{HTML}{5F8FBF}
\definecolor{cvlight}{HTML}{9BBBD9}
\definecolor{cvink}{HTML}{1D2733}
\definecolor{cvmuted}{HTML}{3A4654}
\begin{document}
\begin{tikzpicture}[x=1mm, y=1mm]
  % lap
  \fill[white] (0,0) rectangle (176,250);
  % logó
  \node[anchor=north west, inner sep=0] at (16,232)
    {\includegraphics[width=66mm]{../../assets/logo/pe-cimer-nevjel-fekvo-kek.pdf}};
  % csoportosítás-motívum: három pontfelhő és középpontjaik
  \pgfmathsetseed{2026}
  \foreach \cx/\cy/\col in {128/150/cvblue, 152/122/cvmid, 118/104/cvlight}{
    \foreach \i in {1,...,85}{
      \pgfmathsetmacro{\u}{\cx + 8.5*(rand + rand + rand)}
      \pgfmathsetmacro{\v}{\cy + 7*(rand + rand + rand)}
      \fill[\col, opacity=.55] (\u,\v) circle (.62);
    }
    \draw[cvblue, line width=1.3pt] (\cx-2.4,\cy-2.4) -- (\cx+2.4,\cy+2.4)
      (\cx-2.4,\cy+2.4) -- (\cx+2.4,\cy-2.4);
  }
  % cím
  \node[anchor=base west, inner sep=0, text=cvblue] at (15,176)
    {\cvblack\fontsize{44}{48}\selectfont Adatbányászat};
  \node[anchor=base west, inner sep=0, text=cvink] at (16,162)
    {\cvserif\fontsize{22}{26}\selectfont a hatékonyság eszköze};
  \fill[cvblue] (16.5,150) -- (52,150) -- (50.6,147.6) -- (15.1,147.6) -- cycle;
  \node[anchor=base west, inner sep=0, text=cvmuted] at (16,136)
    {\cvserif\itshape\fontsize{15}{18}\selectfont Elmélet és gyakorlat Pythonban};
  \node[anchor=base west, inner sep=0, text=cvmuted] at (16,128)
    {\cvsans\fontsize{11}{13}\selectfont 2., teljesen átdolgozott kiadás};
  % szerző
  \node[anchor=base west, inner sep=0, text=cvblue] at (16,86)
    {\cvbold\fontsize{19}{22}\selectfont Abonyi János};
  % ferde alsó sáv
  \fill[cvblue] (0,0) -- (176,0) -- (176,46) -- (0,28) -- cycle;
  \node[anchor=base west, inner sep=0, text=white] at (16,13)
    {\cvsans\bfseries\fontsize{10.5}{12}\selectfont Pannon Egyetem\enspace·\enspace Veszprém\enspace·\enspace 2026};
  \node[anchor=base west, inner sep=0, text=white] at (16,6.5)
    {\cvsans\fontsize{8.5}{10}\selectfont nyílt hozzáférésű tankönyvsorozat};
  \node[anchor=base east, inner sep=0, text=white] at (160,13)
    {\cvsans\fontsize{8.5}{10}\selectfont webkönyv\enspace·\enspace PDF\enspace·\enspace EPUB\enspace·\enspace Colab};
  \node[anchor=base east, inner sep=0, text=white] at (160,6.5)
    {\cvsans\fontsize{8.5}{10}\selectfont complexsystemsintelligence.github.io/adatbanyaszat};
\end{tikzpicture}
\end{document}
